\subsection{Модуль lab4}

\subsubsection{KirchhoffTheorem: подсчёт числа остовных деревьев}

\textbf{Вход:} \texttt{Graph const\&}~--- неориентированный граф.

\textbf{Выход:} \texttt{size\_t}~--- число различных остовных деревьев.

\textbf{Описание:} Метод \texttt{countSpanningTrees} строит матрицу
Лапласа $L$ размера $n \times n$ по формуле: $\ell_{ii} = \deg v_i$
на диагонали, $\ell_{ij} = -1$ если ребро $(i,j)$ существует,
иначе $\ell_{ij} = 0$. Матрица строится через вспомогательный шаблон
\texttt{CollectionUtils::makeMatrix}, принимающий лямбду-заполнитель.
Затем удаляется нулевая строка и нулевой столбец~--- получается минор
$L_{00}$ размера $(n-1)\times(n-1)$, определитель которого по теореме
Кирхгофа равен числу остовных деревьев. Определитель вычисляется
рекурсивным разложением по первой строке с построением подматриц:
для каждого элемента $L_{00}[0][p]$ формируется подматрица размера
$(n-2)\times(n-2)$ вычеркиванием нулевой строки и столбца $p$,
после чего рекурсивно вычисляется её определитель. Базовые случаи~---
матрицы $1\times1$ и $2\times2$ вычисляются напрямую. Результат
округляется до ближайшего целого через \texttt{std::llround} и
зажимается снизу нулём через \texttt{std::max}, чтобы исключить
артефакты плавающей точки. Граничные случаи: граф из нуля вершин
возвращает $0$, из одной вершины~--- $1$.
Сложность: $O(n^3)$, пространственная: $O(n^2)$.

\begin{lstlisting}[style=cppstyle, caption={Построение матрицы Лапласа, минора и вычисление определителя (KirchhoffTheorem.cc)}]
size_t KirchhoffTheorem::countSpanningTrees(Graph const& graph);

Matrix laplacian = CollectionUtils::makeMatrix<double>(
    graph.vertexIds(), graph.vertexIds(),
    [&graph](int i, int j) {
        if (i == j)
            return static_cast<double>(graph.degree(i));
        return graph.hasEdge(i, j) ? -1.0 : 0.0;
    });

Matrix minor(n - 1, std::vector<double>(n - 1));
for (size_t i = 1; i < n; ++i)
    for (size_t j = 1; j < n; ++j)
        minor[i - 1][j - 1] = laplacian[i][j];

return static_cast<size_t>(
    std::llround(std::max(0.0, determinant(minor))));

double KirchhoffTheorem::determinant(Matrix const& matrix);

double det = 0.0;
for (size_t p = 0; p < n; ++p) {
    Matrix submatrix(n - 1, std::vector<double>(n - 1));
    for (size_t i = 1; i < n; ++i) {
        size_t col = 0;
        for (size_t j = 0; j < n; ++j)
            if (j != p)
                submatrix[i - 1][col++] = matrix[i][j];
    }
    double cofactor = (p % 2 == 0 ? 1 : -1) * matrix[0][p];
    det += cofactor * determinant(submatrix);
}
\end{lstlisting}

\subsubsection{Boruvka: алгоритм Борувки}

\textbf{Вход:} \texttt{Graph const\&}~--- связный неориентированный
взвешенный граф.

\textbf{Выход:} \texttt{unique\_ptr<Graph>}~--- минимальное остовное
дерево; \texttt{nullptr} если граф ориентирован, пуст или несвязен.

\textbf{Описание:} Метод \texttt{buildMST} реализует алгоритм Борувки.
На этапе инициализации список всех рёбер собирается однократно~---
каждое ребро добавляется ровно один раз благодаря условию \texttt{u < v},
исключающему дублирование для неориентированного графа. Все вершины
добавляются в результирующий граф \texttt{mst} заранее.

Компоненты связности отслеживаются внутренней структурой
\texttt{UnionFind}. Метод \texttt{find} реализует сжатие пути:
при каждом вызове родитель вершины перезаписывается непосредственно
на корень компоненты. Метод \texttt{unite} объединяет две компоненты
по рангу: корень с меньшим рангом подвешивается под корень с большим;
при равных рангах ранг нового корня увеличивается на $1$.

На каждой итерации главного цикла для каждого ребра $(u, v)$
определяются корни компонент \texttt{set\_u} и \texttt{set\_v};
ребро пропускается если они совпадают. Иначе для каждой из двух
компонент обновляется словарь \texttt{cheapest}: кандидат заменяется
только если его вес строго меньше текущего минимума. После обхода
всех рёбер каждое ребро из \texttt{cheapest} перед добавлением
в остов проверяется повторно~--- условие \texttt{find(u) == find(v)}
отсеивает дубли, возникающие когда одно ребро является наилучшим
сразу для двух компонент. Принятое ребро добавляется в \texttt{mst},
компоненты объединяются, счётчик \texttt{components} уменьшается.
Итерации продолжаются пока \texttt{components > 1}. Если за итерацию
ни одного слияния не произошло~--- граф несвязен, возвращается
\texttt{nullptr}.
Сложность: $O(E \log V)$, пространственная: $O(V + E)$.

\begin{lstlisting}[style=cppstyle, caption={Сбор рёбер, выбор дешевейшего ребра для каждой компоненты и слияние (Boruvka.cc)}]
std::unique_ptr<Graph> Boruvka::buildMST(Graph const& graph);
int Boruvka::UnionFind::find(int v);
void Boruvka::UnionFind::unite(int u, int v);

std::vector<Edge> edges;
for (int u : vertices)
    for (auto const& [v, weight] : graph.neighbors(u))
        if (u < v)
            edges.push_back({u, v, weight});

while (components > 1) {
    std::unordered_map<int, Edge> cheapest;

    for (auto const& edge : edges) {
        int set_u = uf.find(edge.from);
        int set_v = uf.find(edge.to);
        if (set_u == set_v) continue;

        auto update_cheapest = [&cheapest](int component,
                                           Edge const& candidate) {
            auto it = cheapest.find(component);
            if (it == cheapest.end()
                    || candidate.weight < it->second.weight)
                cheapest[component] = candidate;
        };
        update_cheapest(set_u, edge);
        update_cheapest(set_v, edge);
    }

    if (cheapest.empty()) return nullptr;

    bool merged_any = false;
    for (auto const& [component, edge] : cheapest) {
        (void)component;
        if (uf.find(edge.from) == uf.find(edge.to)) continue;
        uf.unite(edge.from, edge.to);
        mst->addEdge(edge.from, edge.to, edge.weight);
        --components;
        merged_any = true;
    }
    if (!merged_any) return nullptr;
}
\end{lstlisting}

\subsubsection{PruferCode: кодирование и декодирование кода Прюфера}

\textbf{Вход (encode):} \texttt{Graph const\&}~--- дерево (как правило,
результат алгоритма Борувки).

\textbf{Выход (encode):} \texttt{PruferEncoding}~--- структура с полями:
\texttt{sequence}~--- последовательность Прюфера длины $n-2$;
\texttt{vertices}~--- отсортированный список вершин;
\texttt{edgeWeights}~--- словарь весов рёбер, индексированный
парой $(\min(u,v),\, \max(u,v))$.

\textbf{Описание (encode):} Перед началом кодирования все веса рёбер
сохраняются в \texttt{edgeWeights}~--- это необходимо для последующего
восстановления взвешенного дерева при декодировании, поскольку в ходе
кодирования структура дерева логически разрушается. Степени вершин
хранятся в локальном словаре \texttt{degree}; граф не модифицируется.
На каждом из $n-2$ шагов линейным поиском по отсортированному списку
вершин находится висячая вершина с минимальным номером (степень равна
$1$). Среди её соседей выбирается тот, чья степень строго больше $0$~---
он записывается в \texttt{sequence}. Затем степени обоих концов
уменьшаются на $1$. Для деревьев из $\leq 2$ вершин последовательность
остаётся пустой~--- граничный случай обрабатывается явным ранним
выходом.

\textbf{Вход (decode):} \texttt{PruferEncoding const\&}~--- структура,
полученная при кодировании.

\textbf{Выход (decode):} \texttt{unique\_ptr<Graph>}~--- восстановленное
взвешенное дерево.

\textbf{Описание (decode):} Степени вершин восстанавливаются из кода:
каждая вершина получает начальную степень $1$, затем степень
увеличивается на $1$ за каждое вхождение в \texttt{sequence}. На
каждом шаге линейным поиском по отсортированному списку вершин
находится наименьшая вершина со степенью $1$~--- она не встречается
в ещё не обработанной части кода, поскольку её степень не будет
увеличиваться. Добавляется ребро между найденной вершиной и текущим
элементом \texttt{sequence}; вес берётся из \texttt{edgeWeights} по
ключу $(\min, \max)$, при отсутствии~--- подставляется $1.0$.
Степени обоих концов уменьшаются на $1$. После $n-2$ шагов ровно
две вершины имеют степень $1$~--- они соединяются последним ребром.
Сложность кодирования и декодирования: $O(n^2)$,
пространственная: $O(n)$.

\begin{lstlisting}[style=cppstyle, caption={Кодирование: сохранение весов, поиск листа и заполнение последовательности (PruferCode.cc)}]
PruferCode::PruferEncoding PruferCode::encode(Graph const& tree);

for (int u : vertices)
    for (auto const& [v, weight] : tree.neighbors(u))
        if (u < v)
            result.edgeWeights[{u, v}] = weight;

for (size_t i = 0; i < vertices.size() - 2; ++i) {
    int leaf = -1;
    for (int v : vertices)
        if (degree[v] == 1) { leaf = v; break; }

    for (auto const& [u, weight] : tree.neighbors(leaf)) {
        if (degree[u] > 0) {
            result.sequence.push_back(u);
            degree[leaf]--;
            degree[u]--;
            break;
        }
    }
}
\end{lstlisting}

\begin{lstlisting}[style=cppstyle, caption={Декодирование: восстановление степеней, добавление рёбер и последнее ребро (PruferCode.cc)}]
std::unique_ptr<Graph> PruferCode::decode(
    PruferEncoding const& encoding);

std::unordered_map<int, int> degree;
for (int v : vertices) degree[v] = 1;
for (int v : encoding.sequence) degree[v]++;

for (int v : encoding.sequence) {
    int leaf = -1;
    for (int candidate : vertices)
        if (degree[candidate] == 1) { leaf = candidate; break; }

    tree->addEdge(leaf, v, lookupWeight(leaf, v));
    degree[leaf]--;
    degree[v]--;
}

int v1 = -1, v2 = -1;
for (int candidate : vertices)
    if (degree[candidate] == 1) {
        if (v1 == -1) v1 = candidate;
        else        { v2 = candidate; break; }
    }
if (v1 != -1 && v2 != -1)
    tree->addEdge(v1, v2, lookupWeight(v1, v2));
\end{lstlisting}

\subsubsection{GraphCombinatorics: минимальная раскраска графа}

\textbf{Вход:} \texttt{Graph const\&}~--- произвольный граф
(исходный или МОД~--- по выбору пользователя).

\textbf{Выход:} \texttt{vector<int>}~--- вектор цветов длины $n$,
где элемент с индексом $i$ соответствует $i$-й вершине в
отсортированном порядке; цвета нумеруются с $0$.

\textbf{Описание:} Метод \texttt{findMinColoring} перебирает
число цветов $k$ от $1$ до $n$ включительно. Для каждого $k$
вектор \texttt{colors} сбрасывается в $-1$ и запускается
рекурсивный \texttt{backtrackColoring}. На каждом шаге рекурсии
обрабатывается вершина с индексом \texttt{index} в отсортированном
порядке. Для неё перебираются цвета от $0$ до $k-1$: цвет считается
допустимым, если ни один из уже окрашенных соседей не имеет такого
же цвета~--- соседи ищутся через \texttt{graph.getNeighbors},
их индекс в векторе цветов определяется через \texttt{vertex\_to\_index}.
Если допустимый цвет найден, он присваивается и рекурсия идёт
на следующую вершину; при неудаче цвет сбрасывается в $-1$
и перебирается следующий. Первое $k$, при котором рекурсия
вернула \texttt{true}, является хроматическим числом $\chi(G)$~---
найденная раскраска сохраняется в \texttt{best\_colors} и цикл
прерывается. После завершения правильность раскраски дополнительно
верифицируется методом \texttt{isValidColoring}, обходящим все рёбра
графа.
Сложность в худшем случае: $O(n \cdot k^n)$,
пространственная: $O(n)$.

\begin{lstlisting}[style=cppstyle, caption={Перебор числа цветов, backtracking и проверка допустимости (GraphCombinatorics.cc)}]
VerticesSet GraphCombinatorics::findMinColoring(
    Graph const& graph);

bool GraphCombinatorics::backtrackColoring(
    Graph const& graph, int index,
    VerticesSet const& vertices, VerticesSet& colors,
    int num_colors,
    std::unordered_map<int, int> const& vertex_to_index);

bool GraphCombinatorics::isValidColoring(
    Graph const& graph, VerticesSet const& colors,
    VerticesSet const& vertices,
    std::unordered_map<int, int> const& vertex_to_index) const;

for (int num_colors = 1; num_colors <= static_cast<int>(n); ++num_colors) {
    std::fill(colors.begin(), colors.end(), -1);
    if (backtrackColoring(graph, 0, vertices, colors,
                          num_colors, vertex_to_index)) {
        best_colors = colors;
        break;
    }
}

int vertex = vertices[index];
for (int color = 0; color < num_colors; ++color) {
    bool canColor = true;
    for (int neighbor : graph.getNeighbors(vertex)) {
        int neighbor_idx = vertex_to_index.at(neighbor);
        if (colors[neighbor_idx] == color) {
            canColor = false;
            break;
        }
    }
    if (canColor) {
        colors[index] = color;
        if (backtrackColoring(graph, index + 1, vertices, colors,
                              num_colors, vertex_to_index))
            return true;
        colors[index] = -1;
    }
}
\end{lstlisting}
